% TOP_PROBLEM: {"id": 30006710, "title": "Non-Triviality of the Three-Dimensional Ising Model", "category_id": 16, "category": {"id": 16, "name": "physics", "display_name": "Mathematical Physics", "description": "Problems at the intersection of mathematics and physics.", "slug": "physics", "order_index": 16, "created_at": "2026-07-31T15:26:25.671Z"}, "status": "open"}
% FIRST_POSTED: 2026-09-27
% !TeX program = lualatex
% Compile twice: lualatex 30006710.tex
% Standalone research notebook; original section preserved.
% Research additions: 27 September 2026. No external TeX input or BibTeX file.
\documentclass[11pt,a4paper]{article}
\usepackage[margin=24mm,headheight=14pt]{geometry}
\usepackage{fontspec}
\setmainfont{Latin Modern Roman}
\setsansfont{Latin Modern Sans}
\setmonofont{Latin Modern Mono}
\usepackage{amsmath,amssymb,mathtools}
\usepackage{unicode-math}
\setmathfont{Latin Modern Math}
\usepackage{microtype}
\usepackage{enumitem,xurl,needspace}
\usepackage[dvipsnames]{xcolor}
\usepackage{hyperref}
\hypersetup{unicode=true,colorlinks=true,linkcolor=MidnightBlue,urlcolor=MidnightBlue,
 pdftitle={Research notebook - Problem 30006710},pdfauthor={Open Problems Math},bookmarksnumbered=true}
\usepackage{bookmark}
\usepackage{titlesec}
\titleformat{\section}{\Large\bfseries\raggedright}{\thesection}{0.7em}{}
\usepackage{fancyhdr}
\pagestyle{fancy}\fancyhf{}
\fancyhead[L]{\small\sffamily Open Problems Math | Research notebook}
\fancyhead[R]{\small\sffamily Problem 30006710 | 27 September 2026}
\fancyfoot[C]{\thepage}
\setlength{\parindent}{0pt}
\setlength{\parskip}{5pt plus 1pt}
\setlength{\emergencystretch}{3em}
\setcounter{tocdepth}{1}
\setcounter{secnumdepth}{1}
\setlist[itemize]{leftmargin=3em,itemsep=5pt,topsep=4pt,parsep=0pt}
\allowdisplaybreaks
\numberwithin{equation}{section}
\newcommand{\N}{\mathbb{N}}
\newcommand{\Z}{\mathbb{Z}}
\newcommand{\Q}{\mathbb{Q}}
\newcommand{\R}{\mathbb{R}}
\newcommand{\C}{\mathbb{C}}
\newcommand{\F}{\mathbb{F}}
\newcommand{\A}{\mathbb{A}}
\newcommand{\PP}{\mathbb P}
\newcommand{\E}{\mathbb{E}}
\newcommand{\eps}{\varepsilon}
\newcommand{\dd}{\,\mathrm d}
\newcommand{\sref}[2]{\hyperlink{source:#1:#2}{[S#2]}}
\newcommand{\eref}[2]{\hyperlink{extra:#1:#2}{[E#2]}}
\newif\ifshowcatalogue
\showcataloguetrue
\newcommand{\cataloguescope}[1]{\ifshowcatalogue
 \paragraph{Exact scope in the supplied catalogue.}
 \begin{quote}\small #1\end{quote}\fi}
\begin{document}
\setcounter{section}{0}
\section*{\texorpdfstring{Non-Triviality of the Three-Dimensional Ising Model}{Non-Triviality of the Three-Dimensional Ising Model}}\label{30006710}
\subsection{Definitions and mathematical statement}
At the critical point of the nearest-neighbor three-dimensional Ising model, consider suitably normalized spin fields under vanishing lattice spacing. A centered Gaussian field has all odd moments zero and its even moments equal sums over pairings of two-point correlations. Nontriviality in this record means a scaling-limit field violating this rule; a possible witness is a nonzero truncated four-point correlation
\[
 \langle\Phi_1\Phi_2\Phi_3\Phi_4\rangle
 -\langle\Phi_1\Phi_2\rangle\langle\Phi_3\Phi_4\rangle
 -\langle\Phi_1\Phi_3\rangle\langle\Phi_2\Phi_4\rangle
 -\langle\Phi_1\Phi_4\rangle\langle\Phi_2\Phi_3\rangle.
\]
Conformal invariance is not part of this narrower target. The normalization and mode of convergence must be specified in a rigorous formulation of the scaling limit.
\par\smallskip\textit{Formulation and concept references:} \sref{30006710}{1}.
\cataloguescope{Prove that the critical three-dimensional Ising model is non-trivial, that is, that its scaling limit is not Gaussian and its critical spin correlations do not factorize according to Wick's rule.}
\subsection{Short English statement}
Can one prove that the critical three-dimensional Ising scaling limit is genuinely interacting rather than Gaussian?
% BEGIN RESEARCH_ADDITION ATTEMPT_73
\subsection{Research attempt}

\paragraph{Current frontier.}
The random-current representation turns the truncated four-point function
into a connectivity probability. Aizenman--Duminil-Copin's Section~3.2
provides the exact identity used here \eref{30006710}{1}. Their four-dimensional
Gaussian theorem, and Panis's effective-dimension-at-least-four extension,
do not prove nontriviality of the nearest-neighbour critical model in
three dimensions \eref{30006710}{2}. Numerical exponents or bootstrap consistency
checks in the retained catalogue are not substituted for construction and
moment control of a lattice scaling limit.

\paragraph{Attack route.}
One can seek a lower bound on a normalized four-point cumulant, try to
transfer a continuum interacting theory through universality, or estimate
current intersections directly. Universality would itself need proof for
the lattice model. The route pursued is the third: use a second-moment
intersection estimate to obtain a uniform cumulant gap. The main point is
to keep the normalization, the source sampling, and the passage to a limit
explicit; a positive finite-lattice cumulant is not enough.

\paragraph{Attempt: sampling four sources at the correct scale.}
Work first on a finite connected ferromagnetic Ising graph at zero field.
Let $G(x,y)=\langle\sigma_x\sigma_y\rangle$, choose nonnegative weights
$w_x$ not all zero, and put
\[
 V=\sum_{x,y}w_xw_yG(x,y)>0,
 \qquad X=V^{-1/2}\sum_xw_x\sigma_x.
\]
Then $\langle X\rangle=0$ and $\langle X^2\rangle=1$. Define a probability
law on ordered pairs by
$\nu(x,y)=w_xw_yG(x,y)/V$.
For four distinct sources, the switching identity states
\begin{equation}\label{eq73:current}
 U_4(x,y,z,t)=-2G(x,y)G(z,t)
   \mathbb P^{xy,zt}(x\leftrightarrow z\text{ in }n_1+n_2),
\end{equation}
where the independent currents have respective source sets $\{x,y\}$ and
$\{z,t\}$. A current assigns nonnegative integers to edges with weight
$\prod_e(\beta J_e)^{n_e}/n_e!$, conditioned on the odd-degree source set.
Equation (3.11) of \eref{30006710}{1} is the precise input; no comparison with
independent random walks has been presumed.

Let $\mathcal Q$ be any chosen set of ordered quadruples with distinct
sources. The truncated function is nonpositive, by the same identity for
distinct sources. Coincident-source terms are also nonpositive: for example
$U_4(x,x,z,t)=-2G(x,z)G(x,t)$, using $\sigma_x^2=1$.
Expanding the fourth cumulant and discarding those other nonpositive terms
gives
\begin{equation}\label{eq73:weighted}
 \kappa_4(X)\leq-2\sum_{(x,y,z,t)\in\mathcal Q}
       \nu(x,y)\nu(z,t)
       \mathbb P^{xy,zt}(x\leftrightarrow z).
\end{equation}
This avoids a hidden claim about how repeated-source notation is interpreted
in a current formula. It also makes clear that a lower bound at one fixed
quadruple, carrying vanishing sampling mass, would not suffice.

\paragraph{Attempt: an intersection certificate.}
For each source pair select a path in the positive support of its current,
using a fixed ordering to make the selection measurable. A current with
exactly two distinct sources always connects them: every connected
component has an even number of odd-degree vertices. Let the selected
paths be $\gamma_{xy}$ and $\gamma_{zt}$, and put
$I=|\gamma_{xy}\cap\gamma_{zt}|$.
Whenever $I>0$, the two source pairs belong to one component of $n_1+n_2$.
For each quadruple with $\mathbb E I>0$, Cauchy--Schwarz gives
\[
 (\mathbb E I)^2
   =(\mathbb E[I\mathbf1_{I>0}])^2
   \leq\mathbb E[I^2]\,\mathbb P(I>0).
\]
Consequently the explicit assumptions
\begin{equation}\label{eq73:missing}
 (\nu\otimes\nu)(\mathcal Q)\geq\alpha>0,
 \qquad 0<\mathbb E^{xy,zt}I,
 \qquad \mathbb E^{xy,zt}I^2\leq C(\mathbb E^{xy,zt}I)^2
 \quad((x,y,z,t)\in\mathcal Q)
\end{equation}
imply
\begin{equation}\label{eq73:gap}
 \kappa_4(X)\leq-2\alpha/C.
\end{equation}
All constants must be uniform along the intended vanishing-mesh limit.
A natural choice of $\mathcal Q$ places the four sources in four separated
macroscopic boxes; the first condition then asks that the pair-correlation
sampling measure allocate a positive fraction to these boxes. Neither
that mass bound nor the relative second-moment estimate is asserted here
for critical three-dimensional currents.

This formulation leaves room to replace selected paths by smaller
measurable subsets of the source clusters, provided their intersection
still forces connection. Such a modification might suppress rare large
intersection clusters that inflate $\mathbb E I^2$. It is not enough to
truncate $I$ arbitrarily: one must re-establish a nonzero first moment and
a uniform ratio for the modified witness. For instance replacing $I$ by
$\mathbf1_{I>0}$ makes its ratio exactly $1/\mathbb P(I>0)$, which merely
restates the desired estimate.

\paragraph{Attempt: passing a cumulant gap to a genuine limit.}
To formulate a conditional scaling conclusion, take critical nearest-neighbour
models with lattice spacing $a\downarrow0$, volumes tending to infinity
in a specified order, and weights from one fixed nonnegative compactly
supported test function. Use the actual variance $V_a$ above, not a guessed
critical exponent. Suppose the corresponding scalar variables $X_a$ converge
weakly along a subsequence and
\begin{equation}\label{eq73:ui}
 \sup_a\mathbb E|X_a|^8<\infty.
\end{equation}
The eighth-moment bound implies uniform integrability of powers through
four. Thus the weak limit $X_\infty$ is centred, has variance one, and
satisfies
$\mathbb E X_\infty^4-3\leq-2\alpha/C$ if
\eqref{eq73:missing} holds uniformly. It cannot be Gaussian.
Every limiting random field whose tested marginal is this variable is
therefore non-Gaussian. Scalar tightness follows from variance one, but
this by itself does not construct joint limits for all test functions or
a unique generalized random field; those additional steps remain explicit
requirements of the original scaling-limit target.

\paragraph{Stress test.}
A large $\mathbb E I$ is not a connection lower bound. The elementary law
$I=R$ with probability $1/R$ and $I=0$ otherwise has $\mathbb E I=1$
while $\mathbb P(I>0)\to0$ and $\mathbb E I^2=R$.
This is the precise failure mode the second condition in
\eqref{eq73:missing} must exclude. Upper bounds from a tree diagram go in
the wrong direction to establish it. Moreover, upper bounds on the
spin two-point function do not produce the joint path estimates or a
random-walk Markov property for currents. The source explicitly separates
those difficulties \eref{30006710}{1}, Section~3.2.

Variance normalization also does not justify moment convergence without
\eqref{eq73:ui} or a substitute. Rare tails may disappear in weak
convergence and retain fourth moments. A finite graph computation cannot
supply uniform constants through arbitrarily large critical volumes.
The companion script enumerates a four-cycle at $\tanh\beta=1/2$ using
rational weights, checks a negative normalized cumulant and nonpositive
truncated four-point terms, and is labelled solely as a finite identity
check. This is not three-dimensional critical evidence.

\paragraph{Outcome.}
\textbf{Outcome: Conditional reduction.}

A source-weighted current-intersection criterion is proved to force a
quantitative non-Gaussian marginal in any scaling limit satisfying the
specified moment control. The first- and second-moment roles, repeated
sources, normalization, and weak-limit qualification are all separated.
No uniform three-dimensional intersection estimate has been proved.

The remaining work is to establish \eqref{eq73:missing} at the critical
point with a justified field-level limiting construction and moment control,
or to find another surviving interaction witness. Neither continuum
universality nor numerical critical exponents are assumed. No new
unconditional Ising theorem or priority claim is made.

% Keep the local bibliography together rather than leave a source fragment.
\needspace{170mm}
% END RESEARCH_ADDITION ATTEMPT_73
\subsection{Sources}
\begin{itemize}\raggedright
\item[\textnormal{[S1]}]\hypertarget{source:30006710:1}{} Lectures on the Ising and Potts models on the hypercubic lattice (Hugo Duminil-Copin).
\url{https://arxiv.org/html/1707.00520}.
{\small\textit{Catalogue formulation source.}}
\item[\textnormal{[S2]}]\hypertarget{source:30006710:2}{} Triviality of the scaling limits of critical Ising and phi\textasciicircum{}4 models with effective dimension at least four.
\url{https://arxiv.org/abs/2309.05797}.
{\small\textit{Catalogue research/status reference; not a proof endorsement.}}
\item[\textnormal{[S3]}]\hypertarget{source:30006710:3}{} Lectures on the Ising and Potts models on the hypercubic lattice.
\url{https://www.ihes.fr/~/duminil/publi/2017PIMS.pdf}.
{\small\textit{Catalogue research/status reference; not a proof endorsement.}}
\item[\textnormal{[S4]}]\hypertarget{source:30006710:4}{} Certified three-dimensional Ising critical exponents from harmonized bootstrap series and interval Monte Carlo.
\url{https://link.springer.com/article/10.1140/epjb/s10051-025-01109-8}.
{\small\textit{Catalogue research/status reference; not a proof endorsement.}}
\item[\textnormal{[S5]}]\hypertarget{source:30006710:5}{} Certified Three Dimensional Ising Critical Exponents: author release notes.
\url{https://www.researchgate.net/publication/397730833_Certified_Three_Dimensional_Ising_Critical_Exponents}.
{\small\textit{Catalogue research/status reference; not a proof endorsement.}}
% BEGIN RESEARCH_ADDITION REFERENCES_73
\item[\textnormal{[E1]}]\hypertarget{extra:30006710:1}{} Michael Aizenman and
Hugo Duminil-Copin, \emph{Marginal triviality of the scaling limits of
critical 4D Ising and $\phi^4$ models}, Annals of Mathematics 194 (2021),
163--235; arXiv:1912.07973v4, Lemma 3.3 and Section 3.2, especially (3.11).
\url{https://arxiv.org/html/1912.07973}.
{\small\textit{Inspected 27 September 2026; switching and sign inputs only,
not a transfer of the four-dimensional limit theorem to dimension three.}}
\item[\textnormal{[E2]}]\hypertarget{extra:30006710:2}{} Romain Panis,
\emph{Triviality of the scaling limits of critical Ising and $\varphi^4$
models with effective dimension at least four}, arXiv:2309.05797,
abstract and effective-dimension hypotheses.
\url{https://arxiv.org/abs/2309.05797}.
{\small\textit{Inspected 27 September 2026; the effective-dimension
qualification does not include nearest-neighbour critical three-dimensional
Ising as a proved Gaussian case.}}
% END RESEARCH_ADDITION REFERENCES_73
\end{itemize}

\end{document}
